\section{Worked proof plans and genuinely distinguishing cases}
\label{sec:examples}

\subsection{Delayed equality: both the readable basis and the emitted basis}
For update~\eqref{eq:delayed}, a convenient human basis is
\[
 H=\begin{pmatrix}a-c\\b-d\end{pmatrix},\qquad
 H(F(x,u))=\begin{pmatrix}1&1\\0&1\end{pmatrix}H(x).
\]
Both components vanish initially. This one matrix identity proves $a=c$ after
arbitrarily many steps. It explains why merely checking that the target itself
is conserved is insufficient.

The program retains original pullbacks, so its actual stored basis is
\[
 P=\begin{pmatrix}a-c\\a+b-c-d\end{pmatrix},\qquad
 P(F(x,u))=\begin{pmatrix}0&1\\-1&2\end{pmatrix}P(x).
\]
These bases span the same space, but the second is the one in
\code{results/corpus.json}. Showing the first without this qualification would
misdescribe the experiment.

Now change the last update to $d'=d+u+1$. Start at zero and feed two zero
payloads. The first state is $(0,0,0,1)$, so the original demand still holds.
The second is $(0,0,1,2)$, so $a-c=-1$. The prototype emits exactly this
length-two trace. A one-step test at the initial state would miss the error.

\subsection{Different Fibonacci accumulator presentations}
Compare
\[
 (a,b)'=(b,a+b),\qquad (p,q)'=(p+q,p)
\]
starting at $(a,b,p,q)=(0,1,1,0)$, with output demand $b-p=0$.
The additional relation $a-q=0$ is needed. In the emitted basis
\[
 P_1=b-p,\qquad P_2=a+b-p-q,
\]
the transition matrix is
\[
 \begin{pmatrix}0&1\\1&1\end{pmatrix}.
\]
This is not the additive-polynomial recurrence fragment used to synthesize
power sums: the outputs have Fibonacci growth. The proof comes from a finite
relation between the two state representations, not from guessing a polynomial
closed form for either output.

\subsection{A quadratic identity with arbitrary input magnitudes}
Let
\[
 s'=s+u,\qquad t'=t+2su+u^2,
 \qquad (s_0,t_0)=(0,0).
\]
The goal is $t-s^2=0$. The entire closure has one generator because
\[
 (t+2su+u^2)-(s+u)^2=t-s^2.
\]
There is no bound on input length or magnitude. The update is affine in state
$(s,t)$ with coefficients depending polynomially on the payload; the theorem
of Section~\ref{sec:theory} applies with $D=2$.

The case \code{two_payload_square} replaces $u$ by $uv$. The same one-generator
certificate proves the identity for all lists of rational pairs. This checks
that the implementation really extracts coefficients in multiple fresh inputs
rather than substituting a single test value or confusing state degree with
total degree.

\subsection{A finite-control invariant}
Use two control nodes, even and odd, and transitions
\begin{align*}
 \text{even}\to\text{odd}:&(s,t)\mapsto(s+u,u-t),\\
 \text{odd}\to\text{even}:&(s,t)\mapsto(s-u,u-t).
\end{align*}
Start at the even node with $(s,t)=(0,0)$. Demand $s+t=0$ at the even node and
$s-t=0$ at the odd node. Each destination goal pulls back to its source goal:
\[
 (s+u)-(u-t)=s+t,\qquad
 (s-u)+(u-t)=s-t.
\]
The certificate has a one-dimensional basis at each node. A single unqualified
statement $s=t$ would be wrong on half the control states. This is why control
identity is part of the invariant semantics, not just an optimization key.

\subsection{A simultaneous-update regression}
The update $(a,b)\mapsto(b,a)$ with initial $(1,2)$ preserves $a+b=3$.
Sequentially overwriting the first register before computing the second would
instead produce $(2,2)$ and break the invariant. The search uses simultaneous
substitution; the replay checker evaluates every right-hand side against the
same old state. A dedicated regression case and random substitution comparisons
exercise this distinction.

\subsection{The finite-sample trap}
Consider
\[
 x'=x+\prod_{i=0}^{6}(u-i),\qquad x_0=0,
 \qquad\text{demand }x=0.
\]
Every one-step sample at $u=0,1,\ldots,6$ passes. Repeating only those payloads
for any length would also pass. The identity is nevertheless false.
Coefficient extraction discovers a nonzero constant observation at the initial
state, and provenance reconstruction chooses $u=7$. The actual endpoint is
$x=7!=5040$, so the execution checker rejects the universal equality.

The use of a small grid in the \emph{counterexample constructor} is not the
same mistake as finite-sample verification. The nonzero polynomial has already
been established by exact coefficients; the grid is a proved way of finding a
witness to its nonzeroness.

\subsection{A family with known positive and negative answers}
The synthetic conjugacy family begins with an affine machine
\[
 x'=Ax+Bu+c.
\]
Choose an invertible rational matrix $U$ and offset $v$, and express the same
machine in coordinates $y=Ux+v$:
\[
 y'=U\bigl(AU^{-1}(y-v)+Bu+c\bigr)+v.
\]
The product starts with this relation satisfied. Only its \emph{first}
coordinate is supplied as the output goal, so closure must discover any other
relations needed by the selected matrices. The 24 positive cases use two or
three registers per side and one or two input tags. Their 24 negative partners
add $1$ to the first transformed register update, creating a checked violation.

This generator gives known answers and variation in coefficients, basis
coordinates, and tagged transitions. It is not independent evidence that the
method handles arbitrary real-world Lean functions. Because the family was
constructed from the intended algebraic structure, it is a mechanism test,
not a held-out theorem benchmark.
